\section{Problem Formulation}
\label{sec:problem}

Let $\mathcal{D} = \{d_1, \ldots, d_n\}$ be public litigation documents with metadata and text, and let $P$ be a product specification describing architecture, data provenance, modality, and deployment context.
\sysname{} returns a structured risk report $R^*$ and mitigation candidates $M^*$ that must satisfy an auditable goal predicate $\mathcal{G}$ (\S\ref{sec:goal}).
Report quality is not optimized directly; we evaluate it by severity-label F1 (RQ1), blinded report ratings (RQ2), and the coverage loss $L_\tau$ (RQ5).
The core difficulty is a domain gap: product specifications employ engineering terms (data pipelines, API distribution), whereas litigation documents speak in causes of action, procedural posture, and precedent.
\plre{} (\S\ref{sec:plre}) therefore requires mapping engineering-language attributes onto legal-risk theories while conditioning the support set on authority that remains valid at the query date.

\subsection{Goal Predicate and Stopping Rule}
\label{sec:goal}

\sysname{} treats report generation as goal-conditioned planning.
The predicate is
\begin{equation}
\mathcal{G}(R) = \mathcal{G}_{\mathrm{risk}}(R) \land \mathcal{G}_{\mathrm{evidence}}(R) \land \mathcal{G}_{\mathrm{guide}}(R),
\end{equation}
where $\mathcal{G}_{\mathrm{risk}}$ requires every risk item to carry a severity label, $\mathcal{G}_{\mathrm{evidence}}$ requires at least $k$ supporting documents above relevance threshold $\tau$, and $\mathcal{G}_{\mathrm{guide}}$ requires mitigation text for each High or Medium item.
The controller selects among \textsc{Search}, \textsc{Analyze}, \textsc{Summarize}, \textsc{Map}, \textsc{Verify}, and \textsc{Report}; the deterministic (LLM-free) Goal Checker either terminates, emits a structured critique that closes residual gaps, or escalates unmet conditions.

Let $\Phi(R_t)$ count the unmet atomic sub-conditions of $\mathcal{G}$, and let $\mathcal{T}_c^H$ be the human-labeled High-risk theories of case $c$ that are retrievable---i.e., supported by at least one authority in $\mathcal{D}$ that is valid at $t_q$.
The coverage loss at evidence threshold $\tau$ is
\begin{equation}
L_\tau(c) = \mathbf{1}\!\left[\exists\, T \in \mathcal{T}_c^H : T \notin R^*\right].
\end{equation}

\begin{observation}[Coverage-controlled stopping with escalation]
\label{obs:stop}
Assume the set of risk items under consideration is fixed within a run (the Reporter may rewrite mitigation text but does not introduce new risk atoms that expand $\Phi$).
Under monotone evidence caching and attainable progress, $\Phi(R_t)$ is non-increasing, reaches zero in at most $\Phi(R_0)$ successful progress steps when all atoms are attainable, and otherwise returns the residual atoms at $T_{\max}$.
If $\hat{\tau}$ is selected on exchangeable calibration cases by conformal risk control, then the next-case omission risk satisfies $\mathbb{E}[L_{\hat{\tau}}] \le \alpha$.
\end{observation}

The first clause follows almost directly from the definitions (proof in Appendix~\ref{app:theory}).
The second clause is a standard application of conformal risk control \citep{angelopoulos2023conformalrisk} to $L_\tau$, which is monotone in $\tau$.
It governs evidence coverage only, holds only for exchangeable cases---our temporal split violates it at $\alpha \in \{0.05, 0.10\}$ (\S\ref{sec:experiments})---and its tightness depends on the calibration-set size (150 cases).
